% Generated draft: requires mathematical and source review.
\documentclass[11pt]{article}
\usepackage[margin=25mm]{geometry}
\usepackage{fontspec}
\setmainfont{Latin Modern Roman}
\usepackage{amsmath,amssymb,mathtools,unicode-math}
\setmathfont{Latin Modern Math}
\usepackage{enumitem,needspace,xurl,hyperref}
\setlength{\parindent}{0pt}
\setlength{\parskip}{5pt}
\setlength{\emergencystretch}{3em}
\newcommand{\N}{\mathbb{N}}
\newcommand{\Z}{\mathbb{Z}}
\newcommand{\Q}{\mathbb{Q}}
\newcommand{\R}{\mathbb{R}}
\newcommand{\C}{\mathbb{C}}
\newcommand{\F}{\mathbb{F}}
\newcommand{\A}{\mathbb{A}}
\newcommand{\PP}{\mathbb{P}}
\newcommand{\E}{\mathbb{E}}
\newcommand{\eps}{\varepsilon}
\newcommand{\dd}{\,\mathrm{d}}
\begin{document}
\section*{UnsolvedMath ID: 30003090}
\textbf{Source problem code:} catalogue entry
\subsection{Definitions and mathematical statement}
The supplied question concerns a coordinate \(\beta_i\) of an unspecified parameter \(\beta\), with unknown sparsity denoted by \(k\). It asks: in the moderate-sparsity regime, can one construct a valid confidence interval for \(\beta_i\), and can one construct an adaptive honest confidence interval whose construction does not depend on \(k\)?

A precise formulation requires a statistical experiment and parameter classes that the record does not specify. The following gives a conditional formulation only. Let \(X_n\) have distribution \(P_{\theta,n}\), let \(\Theta_n(k)\) denote a parameter class with sparsity indexed by \(k\), and suppose the target \(\beta_i(\theta)\) is real-valued. These are placeholders for missing source specifications. For a prescribed \(\alpha\in(0,1)\), a measurable interval procedure \(C_n(X_n)\subseteq\R\) has finite-sample honest coverage over \(\Theta_n(k)\) if
\[
\inf_{\theta\in\Theta_n(k)}P_{\theta,n}\!\left\{\beta_i(\theta)\in C_n(X_n)\right\}\geq 1-\alpha.
\]
If honesty is intended over several sparsity classes indexed by a set \(K_n\), replace \(\Theta_n(k)\) by \(\bigcup_{k\in K_n}\Theta_n(k)\). An asymptotic interpretation would instead require
\[
\liminf_{n\to\infty}\inf_{\theta\in\bigcup_{k\in K_n}\Theta_n(k)}P_{\theta,n}\!\left\{\beta_i(\theta)\in C_n(X_n)\right\}\geq 1-\alpha.
\]
The record does not select between these coverage requirements.

Under the intended, but unspecified, experiment and definition of moderate sparsity, the first question asks for a valid interval for \(\beta_i\). The second asks for a single procedure, constructed without the unknown \(k\), that is honest over the intended collection of sparsity classes and adapts to their respective inferential difficulty. The length criterion or other performance requirement defining adaptation is missing. Coverage alone permits the uninformative interval \(\R\), so it does not capture a nontrivial adaptive inference problem. No particular statistical model, sparsity range, length bound, or adaptation theorem is asserted here.
\subsection{Short English statement}
The question asks whether a coordinate \(\beta_i\) admits a useful confidence interval under moderate sparsity, and whether one interval procedure can give honest coverage and adapt to unknown sparsity \(k\). The statistical model, sparsity regime, coverage convention, and criterion for adaptation are missing, so the mathematical formulation remains conditional.
\subsection{Sources}
\begin{itemize}
\item Supplied local UnsolvedMath registry record, permanent ID \(30003090\), title \emph{Adaptive Honest Confidence Intervals Under Moderate Sparsity}, exact source problem code \texttt{catalogue entry}; collection: \emph{Oberwolfach Reports Open Problems}. The record gives proposed year 2016 and names no proposer.
\item Supplied report citation: \emph{Computationally and Efficient Inference for Complex Large-scale Data} (2016), Oberwolfach Reports, DOI 10.4171/owr/2016/16, pp.\ 741--796. Source URL and literature-source link labeled ``OWR report 16/2016'': \url{https://doi.org/10.4171/owr/2016/16}. The background describes Oberwolfach Reports as documenting mathematical workshops and their open-problem discussions; no precise problem location within the report is supplied.
\item Supplied literature-source label ``Cai--Guo'': \url{https://doi.org/10.1214/16-AOS1461}. The supplied literature assessment, checked on \textbf{2026-08-21}, labels the problem partially solved and states that Cai and Guo establish valid minimax intervals across sparsity regimes, with honest adaptation independent of \(k\) possible only in the ultra-sparse region and impossible more generally. This is supplied dated curation derived from the OWR curation, not an independently verified literature assessment. The model assumptions, parameter ranges, and exact theorem statements supporting that assessment require independent verification.
\item The registry labels the statement \texttt{corrected\_verified}, but its verification note specifies only safe notation, typography, and TeX-delimiter cleanup, with the hypotheses, quantifiers, and subquestions preserved. This label does not establish mathematical completeness or correctness.
\end{itemize}
\end{document}
